\subsection{比较关系的积分：II. 强关系}

命题 6. 设 \(\mathbf{f}\) 是一个规则向量函数， \(g\) 是一个 \(\geqslant 0\) 于 \([a, +\infty[\) 上的规则函数。

\(1^{\circ}\) 若积分 \(\int_{a}^{+\infty} g(t) dt\) 收敛，则关系 \(\mathbf{f} \ll g\) （相应地， \(\mathbf{f} \sim \mathbf{c}g\) ，其中 \(\mathbf{c}\) 为常数）蕴含 \(\int_{x}^{+\infty} \mathbf{f}(t) dt \ll \int_{x}^{+\infty} g(t) dt\) （相应地， \(\int_{x}^{+\infty} \mathbf{f}(t) dt \sim \mathbf{c} \int_{x}^{+\infty} g(t) dt\) ）。

\(2^{\circ}\) 若积分 \(\int_{a}^{+\infty} g(t) dt\) 为无穷，则关系 \(\mathbf{f} \ll g\) （相应地， \(\mathbf{f} \sim \mathbf{c}g\) ）蕴含

\[
\int_ {\alpha} ^ {x} \mathbf {f} (t) d t \ll \int_ {\beta} ^ {x} g (t) d t \quad (\text { resp. } \quad \int_ {\alpha} ^ {x} \mathbf {f} (t) d t \sim \mathbf {c} \int_ {\beta} ^ {x} g (t) d t),
\]

对任何 \(\alpha\) 与 \(\beta\) （在 \([a, +\infty[\) 中）。

只需对关系 \(\mathbf{f} \ll g\) 证明本命题，因为若 \(\mathbf{c} \neq 0\) ，关系 \(\mathbf{f} \sim \mathbf{c}g\) 等价于 \(\mathbf{f} - \mathbf{c}g \ll g\) 。

第一部分是中值定理的直接推论，因为若有 \(\|\mathbf{f}(x)\| \leqslant \varepsilon g(x)\) （对 \(x \geqslant x_{0}\) ），则可推出

\[
\left\| \int_ {x} ^ {+ \infty} \mathbf {f} (t) d t \right\| \leqslant \int_ {x} ^ {+ \infty} \| \mathbf {f} (t) \| d t \leqslant \varepsilon \int_ {x} ^ {+ \infty} g (t) d t \quad \text { for } x \geqslant x _ {0}.
\]

其次，设 \(\int_{a}^{+\infty} g(t) dt = +\infty\) 。若有 \(\| \mathbf{f}(x) \| \leqslant \varepsilon g(x)\) （对 \(x \geqslant x_0 \geqslant \max(\alpha, \beta)\) ），则有

\[
\begin{array}{l} \int_ {\alpha} ^ {x} \| \mathbf {f} (t) \| d t = \int_ {\alpha} ^ {x _ {0}} \| \mathbf {f} (t) \| d t + \int_ {x _ {0}} ^ {x} \| \mathbf {f} (t) \| d t \\ \leqslant \int_ {\alpha} ^ {x _ {0}} \| \mathbf {f} (t) \| d t + \varepsilon \int_ {x _ {0}} ^ {x} g (t) d t \\ = \varepsilon \int_ {\beta} ^ {x} g (t) d t + \left(\int_ {\alpha} ^ {x _ {0}} \| \mathbf {f} (t) \| d t - \varepsilon \int_ {\beta} ^ {x _ {0}} g (t) d t\right). \end{array}
\]

现在存在一个 \(x_{1} \geqslant x_{0}\) ，使得对每个 \(x \geqslant x_{1}\)

\[
\left| \int_ {\alpha} ^ {x _ {0}} \| \mathbf {f} (t) \| d t - \varepsilon \int_ {\beta} ^ {x _ {0}} g (t) d t \right| \leqslant \varepsilon \int_ {\beta} ^ {x} g (t) d t
\]

由此，对 \(x \geqslant x_1\)

\[
\left\| \int_ {\alpha} ^ {x} \mathbf {f} (t) d t \right\| \leqslant \int_ {\alpha} ^ {x} \| \mathbf {f} (t) \| d t \leqslant 2 \varepsilon \int_ {\beta} ^ {x} g (t) d t
\]

这就完成了证明，因为 \(\varepsilon > 0\) 是任意的。

换言之，可以把强关系 \(\mathbf{f} \ll g\) 、 \(\mathbf{f} \sim \mathbf{a}g\) 的两项积分（当 \(g\) 在一个区间 \([a, +\infty[\) 上为正时），并且关系在两项的原函数之间保持，只要注意从 \(x\) 积到 \(+\infty\) （若 \(\int_{a}^{+\infty} g(t) dt\) 收敛），而在相反的情形从 \(\alpha\) 积到 \(x\) （对任何 \(\alpha\) ，在 \([a, +\infty[\) 中）。

注意，命题 1 (V, p.~228) 与命题 6 (V, p.~230) 当 \(\mathfrak{F}\) 是区间 \([t, +\infty[\) （其中 \(t > a\) ）在一可数集的余集上的迹滤子的一个基时仍然有效 (cf.~I, p.~15, th. 2)。

例. 1) 把 V, p.~230 的命题 6 应用于关系 \(1 / x \ll x^{\alpha - 1}\) （其中 \(\alpha > 0\) ），重新得到关系 \(\log x \ll x^{\alpha}\) （对每个 \(\alpha > 0\) ），它等价于 III, p.~105 中证明的关系 \(y^{1 / \alpha} \ll e^{y}\) 。

\begin{enumerate}
\def\labelenumi{\arabic{enumi})}
\setcounter{enumi}{1}
\tightlist
\item
  我们有 \(\left(\frac{e^x}{x}\right)' = \frac{e^x}{x}\left(1 - \frac{1}{x}\right) \sim e^x / x\) ；由于 \(e^x / x\) 趋于 \(+\infty\) （随 \(x\) ），由 V, p.~230 的命题 6 推出 \(\int_{1}^{x} \frac{e^t}{t} dt \sim e^x / x\) 。
\end{enumerate}

注. 当不假设 \(g\) 保持 \(\geqslant 0\) 于一个区间 \([a, +\infty[\) 上（或保持 \(\leqslant 0\) 于这样一个区间上），且 \(\int_{a}^{+\infty} g(t) dt\) 不收敛时，关系 \(f \sim g\) 不必然蕴含 \(\int_{a}^{x} f(t) dt \sim \int_{a}^{x} g(t) dt\) ，如下例所示： \(g(x) = \sin x\) 且 \(f(x) = \left(1 + \frac{\sin x}{x}\right) \sin x\) ；这里

\pandocbounded{\includegraphics[keepaspectratio]{images/part002_ecaa53ed.jpg}}

\[
\int_ {n \pi} ^ {(n + 1) \pi} \frac {\sin^ {2} t}{t} d t \geqslant \frac {1}{(n + 1) \pi} \int_ {0} ^ {\pi} \sin^ {2} t d t \geqslant \frac {1}{2} \int_ {n + 1} ^ {n + 2} \frac {d t}{t},
\]

由此得

\[
\int_ {\pi} ^ {n \pi} \frac {\sin^ {2} t}{t} d t \geqslant \frac {1}{2} \int_ {2} ^ {n + 1} \frac {d t}{t}
\]

而积分 \(\int_{1}^{+\infty} dt / t\) 为无穷，尽管 \(\int_{\frac{\pi}{2}}^{x} g(t) dt = -\cos x\) 保持有界（参见 V, p.~260, exerc. 4）。

